\sgaStableTarget{sga6.e07.rem.4-9}\textbf{Remark 4.9.} When @@M0@@ is noetherian and @@M1@@ regular, proceeding as in 4.7 we can describe the graded ring structure of @@M2@@ in terms of the graded rings @@M3@@, @@M4@@, the additive homomorphisms @@M5@@ and @@M6@@ induced by @@M7@@ and @@M8@@, and the element @@M9@@ of @@M10@@. Likewise, the same data describe the homomorphisms @@M11@@ defined from @@M12@@ by passage to the associated gradeds.

\sgaStableTarget{sga6.e07.sec.4.10}\subsection*{4.10.}

Let @@M13@@ be a \emph{regular} scheme having an ample sheaf. We know (IV\qquad) that the Grothendieck groups formed with locally free @@M14@@-Modules of finite type and with coherent @@M15@@-Modules are isomorphic. We use this fact to define a ``topological'' filtration on @@M16@@, taking @@M17@@ to be the subgroup generated by the classes of coherent sheaves whose support has codimension @@M18@@; it is easy to show that @@M19@@ is generated by the classes of the sheaves @@M20@@, where @@M21@@ ranges over the integral closed subschemes of @@M22@@ of codimension @@M23@@(X\qquad).

\sgaStableTarget{sga6.e07.prop.4-11}\textbf{Proposition 4.11.} \uline{Let} @@M24@@ \uline{be a regular scheme having an ample sheaf}. \uline{Then} \uline{for every} @@M25@@:
@@M26@@
@@M27@@ The reverse inclusion, valid without tensoring with @@M28@@, will be proved in X\qquad).

By 4.10, it suffices to show that for an integral closed subscheme @@M29@@ of @@M30@@, of codimension @@M31@@, we have:
@@M32@@

For this, we proceed by noetherian induction on the integral closed subschemes of @@M33@@. Suppose therefore that the property is proved for the integral closed subschemes of @@M34@@, and let us prove it for @@M35@@. Let @@M36@@ be the Ideal of @@M37@@
